9.0 WALKA

Ogólna zasada:

Walka może (ale nie musi) mieć miejsce wtedy, gdy oddział natknie się na potwora lub grupę potworów.
Składa się ona z szeregu rund.
Rundy są to powtarzalne bloki określonych czynności (faz), które ma do wykonania każda z walczących stron.
Rozpoczęta walka nie może zostać przerwana do chwili śmierci jednej ze stron.

Sposób postępowania:

Przed rozpoczęciem walki gracze muszą ustalić liczbę oraz rodzaj napotkanych potworów według zasad, podanych w punkcie 8.0.

Następnie gracze ustalają siłę każdego z napotkanych potworów osobno.
Określa ją kod podany na żetonach od np. 1+1 (patrz 3.2)

Kolejną czynnością jest ustalenie szyku grupy potworów.

Po wykonaniu tych wstępnych czynności walka może się rozpocząć (patrz 9.2).

Zasady szczegółowe:

[9.1] Przed rozpoczęciem walki ustalany jest szyk bojowy grupy potworów.
Oddział rozpoczyna walkę w takim szyku, w jakim znajdował się w chwili spotkania.
Jeżeli spotkano więcej niż 3 potwory należy w pierwszym szeregu ustawić trzy najsilniejsze (o najwyższym współczynniku siła).
Pozostałe są ustawiane w drugim i ewentualnie następnych szeregach.

Gracze mogą wybrać, który potwór będzie się to się na tylko jednego potwora decyduje jego sąsiad, naprzeciw każdej postaci z pierwszego szeregu szyku oddziału.
Jeżeli napotknięto tylko jednego potwora ustawia się go naprzeciw postaci, stojącej na środku pierwszego szeregu szyku oddziału.
Jeżeli są dwa potwory ustawia się je naprzeciw postaci ci skrajnych.

[9.2] Po obliczeniu siły każdego z potworów oraz ustaleniu szyku rozpoczyna się pierwsza runda, która (podobnie jak i każda następna) składa się z szeregu faz (czynności):

1.
Atak oddziału – każdy członek oddziału może zaatakować potwora, stojącego w pierwszym szeregu szyku grupy potworów.
Postacie z pierwszego szeregu szyku oddziału mogą atakować tylko przy pomocy broni (z wyjątkiem łuków i noży), zaś postacie, stojące w drugim szeregu szyku mogą strzelać z łuku, rzucać nożami lub rzucać czary.
Wynik każdego indywidualnego ataku jest natychmiast wprowadzany w życie.

2.
Atak potworów – każdy potwór, znajdując cy się w pierwszym lub drugim szeregach szyku atakuje postać z pierwszego szeregu szyku oddziału.
Potwory, stojące w pierwszym szeregu mogą atakować zarówno przy pomocy broni (zwykłe naturalne, jak pazury, kły, itp.), jak i przy pomocy czarów (jeśli są do tego zdolne).
Potwory z drugiego szeregu szyku mogą rzucać czary (jeśli są do tego zdolne).
Wynik każdego indywidualnego ataku jest natychmiast wprowadzany w życie.

3.
Atak Czarnych Wrót – ta faza znajduje zastosowanie tylko wtedy, gdy oddział jest w tym samym segmencie, co Czarne Wrota (czyli wtedy, gdy zostały one odnalezione).
Czarne Wrota, bez względu na miejsce w szyku potworów mogą użyć w ciągu tej fazy 3 czarów ,,Wybuch'' (patrz 16.0).

4.
Reorganizacja szyku oddziału – gracze mogą z każdego szeregu szyku oddziału przesunąć jedną postać do innego szeregu.
Na końcu tej fazy w żadnym szeregu nie może być więcej niż trzy postacie (łącznie z potworami, znajdującymi się pod wpływem uroku).

5.
Reorganizacja szyku grupy potworów – szyk potworów musi być zreorganizowany w ten sposób, aby w pierwszym szeregu znajdowały się trzy potwory (jeśli to możliwe).

Wykonanie wszystkich powyższych faz stanowi jedną rundę walki.
Następna runda (jeśli jest jeszcze z kim walczyć) rozpoczyna się od początku, od fazy 1.

[9.3] Cel każdego indywidualnego ataku zależy od wzajemnej pozycji postaci i potworów.
Jeżeli trzy postacie walczą z trzema potworami każda z nich, jak również każdy z potworów, może zaatakować przeciwnika, stojącego go na wprost:

OBRAZEK
[Diagram: 3 rzędy postaci (1, 2, 3) i potwory (A, B, C) – A zaatakuje 1, B zaatakuje 2, C zaatakuje 3 i odwrotnie]

Jeżeli w jednym lub obu wrogich szykach jest mniej niż trzech walczących, każdy z nich musi zaatakować przeciwnika stojącego naprzeciw niego, chyba że w takiej samej odległości od atakującego znajduje się więcej niż jeden wróg.
W tym przypadku cel ataku potworów (tylko) ustalany jest losowo, za pomocą rzutu kostką:

OBRAZEK
[Diagram: potwory 1, 2, 3 vs postacie A, B]

Jeżeli A i B są potworami to cel ich ataku ustalany jest przez rzut kostką dla każdego z nich osobno (po jednym rzucie dla każdego z potworów).
Jeżeli wyrzucono 1, 2 lub 3 oczka, to celem ataku A jest 1, zaś celem B jest 2.
Jeżeli wyrzucono 4, 5 lub 6 oczek, to celem ataku A jest 2, zaś B – 3.

Jeżeli potworami są 1, 2 i 3, to 1 zaatakuje A, 3 zaatakuje B, zaś 2 zaatakuje A (1, 2 lub 3 oczka) lub B (4, 5 lub 6 oczek).
Postacie wybierają cel ataku według własnej woli, więc jeśli w przykładzie powyżej A i B są postaciami, to A zaatakuje 1 lub 2, zaś B zaatakuje 2 lub 3.
Podobnie w przykładzie poniżej – postać A może zaatakować któregokolwiek z trzech przeciwników.

OBRAZEK
[Diagram: postać A vs potwory 1, 2, 3]

Jeżeli A jest potworem cel jego ataku zależny będzie od rzutu kostką i będzie nim 1 przy wyrzuconych 1 lub 2 oczkach, 2 – przy 3 lub 4 oczkach, 3 – przy 5 lub 6 oczkach.

Powyższe przykłady nie wyczerpują oczywiście wszystkich możliwych kombinacji.
Stanowią one jedynie wskazówkę, w jaki sposób gracze powinni decydować o celach ataków w innych przypadkach.

Postać atakująca z drugiego szeregu szyku oddziału może wybrać dowolny cel z pierwszego szeregu szyku potworów.
Potwór, atakujący z drugiego szeregu szyku ustala cel ataku tylko przy pomocy rzutu kostką.

[9.4] Postacie mogą atakować tylko podczas fazy 1 – atak oddziału, zaś potwory tylko podczas fazy 2 – atak potworów.
Przed rozpoczęciem walki każda postać powinna zadeklarować jakiej broni użyje.
W czasie walki broń może zostać zmieniona, jednak wymaga to opuszczenia jednej własnej fazy ataku.
Obie walczące strony obowiązują następujące ogólne zasady:

1.
Atakować można tylko raz w ciągu jednej rundy walki (wyjątek – Czarne Wrota = 3 ataki).

2.
W ciągu jednej rundy można zaatakować tylko jednego przeciwnika (wyjątek – Czarne Wrota).

3.
Każdy, zarówno postać, jak i potwór, może zostać zaatakowany przez dowolną liczbę wrogów (oczywiście nie więcej niż sześciu – trzech z pierwszego szeregu wrogiego szyku i trzech z drugiego).

4.
Rany, odniesione w czasie walki muszą być natychmiast zaznaczone na karcie cech postaci lub, w przypadku potworów, na dowolnej dogodnej kartce.

[9.5] Wynik ataku, czyli zadana przeciwnikowi liczba ran, odnajdowany jest w tabeli 9.9 ,,Rezultat walki'' na przecięciu kolumny, odpowiadającej używanemu rodzajowi broni oraz rzędu, odpowiadającemu zmodyfikowanej liczbie oczek, wyrzuconych 1K6.

Przykład: Bohater Almuryk atakuje mieczem – gracz rzuca 1K6 i okazuje się, że wyrzucił 4 oczka.
Dodaje do nich 2 (biegłość Almuryka jest równa ,,+2 miecz'') oraz 3 (wielkość wsp. zręczności), a więc zmodyfikowana liczba oczek wynosi 9.
Następnie gracz w tabeli 9.9, w kolumnie ,,miecz'' sprawdza rząd, odpowiadający 9 oczkom.
Liczba na przecięciu wynosi 1, więc Almuryk w tym ataku zadał wrogowi 1 ranę.

Jednym nożem można rzucić tylko raz w czasie całej walki.
Po zwycięstwie można oczywiście ten nóż wziąć z powrotem.
Z łuku można strzelać dowolną liczbę razy w ciągu walki, ale tylko raz podczas jednej rundy i tylko z drugiego szeregu.

[9.6] Potwory, które nie posiadają określonej broni w celu ustalenia rezultatu ataku używają kolumny ,,potwory'', zaś te, które mają określoną broń (patrz 8.2) posługują się kolumną odpowiednią do jej rodzaju.
Na żetonach reprezentujących potwory, które używają broni innej niż naturalna wydrukowany jest obok współczynnika zręczności litera ,,w''.
Rodzaj broni określony jest w tabeli 8.2 ,,Charakterystyka potworów''.

Podobnie jak w przypadku postaci współ-

czynnik zręczności potwora dodawany jest do wyrzuconej przy ataku liczby oczek.

[9.7] Wynikiem ataku może być zadanie wrogowi określonej liczby ran (patrz 9.5).

Każda rana, odniesiona przez postać musi być zanotowana na jej karcie cech, zaś rany, zadane potworowi należy zaznaczać na kartce papieru, na której wcześniej została zanotowana siła tego potwora.

Śmierć w walce następuje wtedy, gdy liczba odniesionych ran jest równa współczynnikowi siła lub od niego większa.
Dotyczy to zarówno postaci, jak i potworów.

[9.8] W wyniku odniesionych ran szyk obu walczących stron może wymagać reorganizacji.
Dokonuje się ona w czasie odpowiedniej fazy (patrz 9.2).

Reorganizacja szyku oddziału:

1.
Postać może zostać przesunięta na inną pozycję w tym samym szeregu szyku, jeżeli nie zmienia to wzajemnego układu postaci w tym szeregu, tzn. postacie, stojące w tym samym szeregu nie mogą się zamieniać miejscami, chociaż ich odległość od siebie może się zmieniać.
Dotyczy to szeregów, w których jest mniej niż trzy postacie.

2.
Postać może zostać przesunięta z jednego szeregu do drugiego, jednak nie więcej niż o jeden szereg w czasie jednej rundy.

3.
Podczas zmiany szeregu postać może równocześnie zmienić pozycję w szeregu, np. ze skrajnej na środkową lub odwrotnie.

4.
Jeżeli w pierwszym szeregu została przy życiu tylko jedna postać, tej jedynej ruch, jaki może (i musi) być wykonany polega na przesunięciu jej z drugiego szeregu do pierwszego.

5.
Jeżeli w pierwszym szeregu zostały dwie postacie żadna z nich nie może zostać wycofana do drugiego szeregu.

Reorganizacja szyku grupy potworów:

1.
Jeżeli w pierwszym szeregu jest mniej niż trzy potwory jeden lub dwa muszą zostać przesunięte z drugiego szeregu, by uzupełnić lukę.

2.
Jeżeli w pierwszym szeregu są trzy potwory jeden z potworów z dalszego szeregu może zostać przesunięty do przodu.

3.
Potwór może zmieniać pozycję w ramach tego samego szeregu, jednak jeżeli jest w nim dwie wolne pozycje rzut kostką decyduje, którą z nich potwór zajmie.

4.
Potwory nie mogą zamieniać się miejscami w tym samym szeregu, ani przesuwać się do tyłu.

[9.8] Tabela ,,Rezultat walki'' – patrz plansza z tabelami